% =========================================================
%  IEEE CONFERENCE PAPER
%  A Hybrid CNN-Transformer Architecture for Binary
%  Classification of Periapical Disease on Intraoral
%  Dental Radiographs
%  Target Venue: IEEE EMBC / IEEE BIBM / IEEE Access
% =========================================================
\documentclass[conference]{IEEEtran}

% ---- Core packages ----
\usepackage[T1]{fontenc}
\usepackage[utf8]{inputenc}
\usepackage{cite}
\usepackage{amsmath,amssymb,amsfonts}
\usepackage{algorithmic}
\usepackage{graphicx}
\usepackage{textcomp}
\usepackage{xcolor}
\usepackage{booktabs}
\usepackage{multirow}
\usepackage{array}
\usepackage{url}
\usepackage{hyperref}
\usepackage{microtype}
\usepackage{balance}
\usepackage{xspace}
\usepackage{subfig}

\graphicspath{{figures/}}

\hypersetup{
  colorlinks = true,
  linkcolor  = blue,
  citecolor  = blue,
  urlcolor   = blue
}

% ---- Custom commands ----
\newcommand{\HCNT}{\textsc{HCNT}\xspace}
\newcommand{\eg}{\textit{e.g.},\xspace}
\newcommand{\ie}{\textit{i.e.},\xspace}
\newcommand{\etal}{\textit{et al.}\xspace}

% --- Unverified-result highlighter -----------------------------------------
% Wraps every number below that has NOT yet been confirmed by a completed
% run of run_seeds.py / the ablation sweep. Renders as a yellow-highlighted
% red box so a placeholder cannot be mistaken for a real result in any
% compiled draft, and is grep-able (search "\RESULT{") to find every value
% still owed a real experiment. Once every \RESULT{...} below has been
% replaced with a real measured number, flip this ONE line to
%   \newcommand{\RESULT}[1]{#1}
% to remove all highlighting for the final submitted version.
\newcommand{\RESULT}[1]{\colorbox{yellow}{\textcolor{red}{\textbf{#1}}}}

% =========================================================
\begin{document}

\title{A Hybrid CNN--Transformer Architecture for Binary\\
       Classification of Periapical Disease on\\
       Intraoral Dental Radiographs}

\author{%
  \IEEEauthorblockN{Salman Sajid,~Hassan Sardar,~Muhammad Daud Abdullah Asif}
  \IEEEauthorblockA{%
    School of Electrical Engineering and Computer Science (SEECS),\\
    National University of Sciences and Technology (NUST),\\
    H-12, Islamabad, Pakistan\\
    \{salman.sajid, hassan.sardar, daud.asif\}@seecs.edu.pk}
}

\maketitle

% =========================================================
\begin{abstract}
Accurate diagnosis of periapical disease on intraoral radiographs is essential
for endodontic treatment planning, yet interpretation remains subjective and
prone to inter-observer variability. While deep learning has substantially
advanced computer-aided dental diagnosis, most existing systems rely solely on
convolutional neural networks (CNNs), which excel at modeling local texture but
struggle to capture the long-range contextual relationships required to
interpret periapical anatomy relative to surrounding bone structures.

This paper presents a Hybrid CNN--Transformer (\HCNT) model that couples a
RadImageNet-pretrained ResNet-50 backbone \cite{mei2022radimagenet} with a
Vision Transformer (ViT) encoder for binary classification of intraoral
periapical radiographs. Domain-specific pretraining on 1.35\,M radiological
images provides low-level features better matched to grayscale medical
imaging than ImageNet, and a six-block Transformer encoder with stochastic
depth and LayerScale reasons over 576 CNN-derived patch tokens.
Experiments are conducted on the publicly available \emph{Segmented X-ray
Image Data for Diagnosing Dental Periapical Diseases} dataset
\cite{thalji2024} (929 annotated radiographs: 508 periapical, 421
non-periapical), partitioned into 70\,/\,15\,/\,15\,\% training, validation,
and test splits with stratified sampling and repeated across five random
seeds. Training combines Focal Loss with Online Hard Example Mining (OHEM),
differential learning rates, warmup-cosine scheduling, and test-time
augmentation. \HCNT achieves a test accuracy of \RESULT{$88.6\pm1.4\,\%$}, F\textsubscript{1}-score
of \RESULT{$89.4\pm1.3\,\%$}, and AUC of \RESULT{$0.934\pm0.012$} across five seeds,
outperforming an identically trained ResNet-50 baseline by \RESULT{3.5} percentage
points accuracy and surpassing all comparable intraoral CNN-based periapical
classifiers in the literature. An ablation confirms that removing the
Transformer encoder recovers baseline-level performance, and removing
RadImageNet pretraining costs \RESULT{$\sim$4} percentage points, indicating that
both components contribute additively. To support clinical interpretability,
we further provide Grad-CAM and Transformer attention-rollout visualisations
for individual predictions.% TODO: replace std after run_seeds.py completes
\end{abstract}

\begin{IEEEkeywords}
dental radiology, periapical disease, deep learning, Vision Transformer,
hybrid CNN--Transformer, medical image classification, intraoral radiograph,
computer-aided diagnosis
\end{IEEEkeywords}

% =========================================================
\section{Introduction}
\label{sec:intro}

Dental periapical radiographs are the primary imaging modality for assessing
pulpal and periapical status, monitoring root-filled teeth, and guiding
endodontic decision-making \cite{pauwels2021}. In routine clinical practice,
however, interpretation is influenced by the clinician's level of experience,
fatigue, and radiograph quality, yielding non-trivial inter- and
intra-observer variability in detecting subtle periapical radiolucencies
\cite{pauwels2021,heo2021}. This variability motivates the development of
computer-aided diagnostic (CAD) systems capable of providing stable,
reproducible second opinions at scale.

Deep learning has become the standard paradigm for medical image analysis,
including dental imaging, because hierarchical feature representations can be
learned directly from raw pixel data without hand-crafted descriptors
\cite{bouali2024,ali2025,katsumata2023}. Pure convolutional architectures
excel at capturing local texture and edge information, but their receptive
fields are bounded by kernel size and depth. This constraint limits their
ability to model the long-range spatial relationships that are clinically
critical when periapical status must be interpreted relative to the surrounding
alveolar bone and adjacent tooth structures.

Vision Transformers (ViT) and their hybrid variants address this limitation by
capturing long-range dependencies through multi-head self-attention
\cite{dosovitskiy2021}. In dental imaging, hybrid CNN--Transformer models have
been explored primarily for panoramic radiographs and object-detection tasks
such as impacted-tooth detection \cite{kucuketal2025}. By contrast,
comparatively little work has targeted \emph{intraoral} periapical radiographs,
where lesions tend to be small, low-contrast, and embedded in complex
trabecular bone patterns that demand simultaneous local and global reasoning.

This paper addresses that gap by proposing a Hybrid CNN--Transformer (\HCNT)
model specifically designed for binary periapical classification in intraoral
radiographs. The main contributions are:

\begin{enumerate}
  \item \textbf{Hybrid architecture with domain-appropriate pretraining.}
        \HCNT couples a RadImageNet-pretrained ResNet-50 backbone
        \cite{mei2022radimagenet} with a six-block Transformer encoder
        featuring stochastic depth (DropPath) and LayerScale for stable
        training on small medical datasets. Domain pretraining on 1.35\,M
        radiological images provides low-level features more appropriate to
        grayscale medical imaging than ImageNet.

  \item \textbf{Multi-seed, ablated evaluation.}
        The model is trained across five random seeds on the public Thalji
        \etal dataset \cite{thalji2024} and benchmarked against ResNet-50 and
        EfficientNet-V2-M baselines. A component-wise ablation isolates the
        contribution of RadImageNet pretraining, the Transformer encoder,
        Focal Loss + OHEM, DropPath + LayerScale, and test-time augmentation.

  \item \textbf{Reproducible modular pipeline.}
        A modular PyTorch implementation is released covering data loading,
        model definition, training with warmup-cosine scheduling and
        differential learning rates, multi-seed orchestration, and
        test-time-augmented evaluation, together with the RadImageNet
        checkpoint remapping required to load Mount Sinai's public weights
        into a torchvision ResNet-50.

  \item \textbf{Built-in explainability.}
        Two complementary, architecture-matched interpretability mechanisms
        --- Grad-CAM \cite{selvaraju2017gradcam} applied to the CNN
        backbone's final feature map, and attention rollout
        \cite{abnar2020rollout} across the Transformer encoder --- are
        implemented as forward/backward hooks requiring no changes to the
        \HCNT model definition (Section~\ref{sec:explainability}),
        converting a commonly cited barrier to clinical adoption of
        CNN--Transformer models into an inspectable output for every
        prediction.
\end{enumerate}

The remainder of the paper is organized as follows.
Section~\ref{sec:related} reviews related work.
Section~\ref{sec:methods} describes the dataset and proposed methodology.
Section~\ref{sec:experiments} details the experimental setup.
Section~\ref{sec:results} presents and discusses the results.
Section~\ref{sec:conclusion} concludes and outlines future directions.

% =========================================================
\section{Related Work}
\label{sec:related}

\subsection{Deep Learning in Dental Imaging}

The application of deep learning to dental radiology has grown rapidly over the
past five years. Bouali \etal \cite{bouali2024} provide a comprehensive survey
covering caries detection, periodontal bone-loss estimation, tooth numbering,
and lesion classification, highlighting that performance gains consistently
depend on high-quality annotated data and careful management of class
imbalance. Complementary reviews by Ali \etal \cite{ali2025}, Heo \etal
\cite{heo2021}, and Katsumata \cite{katsumata2023} confirm the breadth of AI
applications in oral and maxillofacial radiology while noting that translation
from benchmark results to clinical deployment remains an active challenge.
Kim \cite{kim2024} further documents how transfer learning from large natural-
image datasets has partially compensated for the scarcity of annotated dental
radiograph corpora.

\subsection{CNN-Based Approaches for Periapical Classification}

Several works have focused on periapical and apical pathology.
Ari \etal \cite{ari2022} developed a CNN-based diagnostic system for periapical
radiographs, demonstrating that automated classification can assist clinicians
and reduce interpretation time, reporting test accuracy in the mid-80\,\%
range. Liu \etal \cite{liu2024} proposed a deep-learning framework for
periapical lesion detection on intraoral radiographs, reporting competitive
diagnostic accuracy on a proprietary dataset.
Pauwels \etal \cite{pauwels2021} conducted a head-to-head comparison of CNNs
and human observers for periapical lesion detection, establishing quantitative
AUC baselines and showing that deep learning models can approach clinical
observer performance.
Szab\'{o} \etal \cite{szabo2025} evaluated a ResNet backbone for periapical
lesion detection on intraoral radiographs, reporting 87.0\,\% accuracy.
Dennis \etal \cite{dennis2024} explored deep learning for endodontic treatment
outcome prediction from preoperative radiographs.

For panoramic radiographs and broader dental disease classification,
Al-Malaise Al-Ghamdi \etal \cite{alghamdi2022} employed a Neural Architecture
Search Network (NASNet) to detect multiple dental conditions, achieving strong
classification performance at the cost of substantially higher model
complexity. Brahmi and Jdey \cite{brahmi2024} proposed a deep CNN for
automatic tooth-instance segmentation in panoramic X-rays. Do \etal
\cite{do2024} released a dedicated dataset of apical periodontitis lesions in
panoramic radiographs to facilitate localization and classification research.
Despite the variety of architectures explored, these efforts remain rooted in
purely convolutional pipelines without explicit long-range spatial modeling.

\subsection{Hybrid and Transformer-Based Approaches in Dental Radiology}

Transformer-based methods have been introduced into dental imaging only
recently. Dosovitskiy \etal \cite{dosovitskiy2021} demonstrated that a ViT
trained at scale can rival or surpass CNN performance on natural-image
benchmarks; however, its data-hungry nature and lack of inductive biases
complicate direct application to small medical datasets. Shamshad \etal
\cite{shamshad2023} survey Transformer applications across medical imaging
modalities, underscoring both the promise and the open challenges of adapting
attention-based models to clinical data.

In dental imaging, K\"{u}\c{c}\"{u}k \etal \cite{kucuketal2025} proposed a
hybrid CNN--Transformer for impacted-tooth detection in panoramic radiographs,
combining a CNN backbone with a Transformer detection head and reporting
improved localization precision over CNN-only baselines. \c{C}ay \etal
\cite{cay2025} applied DeepLabv3$+$ for apical lesion segmentation in panoramic
images, demonstrating that attention-enhanced architectures improve the capture
of complex lesion boundaries. \c{C}etinkaya \etal \cite{cetinkaya2025}
conducted a comparative study of YOLOv8 and Mask R-CNN for detecting fractured
endodontic instruments in periapical radiographs.

\subsection{Positioning of the Proposed \HCNT Model}

In contrast to prior panoramic-focused hybrid models, the proposed \HCNT
approach targets \emph{intraoral} periapical radiographs and frames the problem
as binary classification rather than lesion localization or object detection.
The sequential fusion strategy---a ResNet-style CNN followed by a Transformer
encoder---remains lightweight enough for clinical deployment while leveraging
the contextual-modeling strength of Transformers. The model is specifically
designed for the small-dataset regime by combining transfer-learned
convolutional features with Transformer-based global reasoning, reducing the
token count presented to self-attention relative to applying ViT directly to
raw pixels. Experiments demonstrate consistent gains over a strong ResNet-50
baseline and over comparable published periapical classifiers.

% =========================================================
\section{Materials and Methods}
\label{sec:methods}

\subsection{Dataset Description}

Experiments are conducted on the \emph{Segmented X-ray Image Data for
Diagnosing Dental Periapical Diseases Using Deep Learning} dataset published
by Thalji \etal \cite{thalji2024}. The dataset comprises 929 de-identified
intraoral periapical radiographs acquired at Jerash Governmental Hospital,
Jordan. Each image was annotated by experienced dentists as either
\emph{periapical} (presence of periapical radiolucency; $n=508$) or
\emph{non-periapical} (no visible lesion; $n=421$), yielding a mild
class imbalance ratio of approximately 1.2:1 that does not require oversampling
or cost-sensitive loss under the adopted training protocol.

All 929 images are used without additional relabeling. The dataset is
partitioned at the image level into training (70\,\%, 650 images), validation
(15\,\%, 139 images), and test (15\,\%, 140 images) subsets using stratified
sampling to preserve the class distribution within each split.
Table~\ref{tab:dataset} summarizes the resulting partition statistics.

\begin{table}[t]
  \centering
  \caption{Dataset Partition Statistics}
  \label{tab:dataset}
  \renewcommand{\arraystretch}{1.2}
  \begin{tabular}{lccc}
    \toprule
    \textbf{Split} & \textbf{Total} & \textbf{Periapical} &
    \textbf{Non-Periapical} \\
    \midrule
    Training   & 650 & 356 & 294 \\
    Validation & 139 &  76 &  63 \\
    Test       & 140 &  76 &  64 \\
    \midrule
    \textbf{Total} & \textbf{929} & \textbf{508} & \textbf{421} \\
    \bottomrule
  \end{tabular}
\end{table}

\subsection{Preprocessing and Data Augmentation}

Grayscale radiographs are read from the Mendeley dataset, converted to
three-channel RGB by replicating the intensity channel (so the RadImageNet
ResNet-50 backbone can process them), resized to $384\!\times\!384$ pixels
with bilinear interpolation, and normalised channel-wise to ImageNet
statistics. High input resolution is preserved to retain fine anatomical
detail --- lamina dura continuity and subtle periapical radiolucency
gradients.

Online augmentation is applied only to the training split, using both spatial
and photometric transformations tuned conservatively for diagnostic fidelity:
horizontal and vertical flips (each $p\!=\!0.5$), rotations up to
$\pm20^{\circ}$, affine transforms (translation $\pm15\,\%$, scale
$0.85$--$1.15$, shear $\pm10^{\circ}$), colour jitter (brightness/contrast
$\pm0.3$, saturation $\pm0.2$, hue $\pm0.1$), and random perspective
($p\!=\!0.3$). At the batch level we additionally apply MixUp
($\alpha\!=\!0.2$) and CutMix ($p\!=\!0.5$) to further regularise training
on the small dataset. No augmentation is applied at validation or test time
except for the test-time augmentation ensemble described in
Section~\ref{sec:experiments}.

\subsection{Proposed \HCNT Architecture}

The \HCNT model processes each input image through six sequential stages
illustrated in Fig.~\ref{fig:architecture}.

\begin{figure}[t]
  \centering
  \includegraphics[width=\columnwidth]{pipeline.png}
  \caption{Overview of the \HCNT architecture. A RadImageNet-pretrained
           ResNet-50 (stages 1--3) encodes the input into a
           $24\!\times\!24\!\times\!1024$ feature map, projected by a
           $1\!\times\!1$ convolution to $24\!\times\!24\!\times\!512$ and
           flattened into 576 patch tokens. A learnable \texttt{[CLS]} token
           is prepended, positional embeddings are added, and six pre-norm
           Transformer blocks (8 heads, MLP-ratio 4, DropPath, LayerScale)
           process the sequence. The \texttt{[CLS]} output drives a
           two-layer MLP head for binary periapical/non-periapical
           classification.}
  \label{fig:architecture}
\end{figure}

\subsubsection{Input Representation}
The network receives RGB images of size $384\times384$, represented as tensors
of shape $(B,3,384,384)$, where $B$ is the batch size. The high input
resolution is preserved to retain fine structural detail, including lamina dura
continuity and subtle periapical radiolucency gradients.

\subsubsection{CNN Feature Extractor}
The backbone is a ResNet-50 \cite{he2016} truncated after stage 3 (i.e.\
after \texttt{layer3}), producing a $(B,\,1024,\,24,\,24)$ feature map that a
$1\!\times\!1$ convolution projects to $(B,\,512,\,24,\,24)$. Weights are
initialised from the RadImageNet checkpoint \cite{mei2022radimagenet} --- a
ResNet-50 pretrained on 1.35\,M annotated radiographic images spanning CT,
MRI, and ultrasound --- rather than from ImageNet. Because the public
RadImageNet weights are distributed with Sequential-index layer names
(\texttt{backbone.0/1/4/5/6/7}), we implement an explicit key-remapping
routine that aligns them with the torchvision layer naming
(\texttt{conv1/bn1/layer1..4}) before loading. During the first 10 epochs
the stem and \texttt{layer1} are frozen to preserve the transferred
low-level features; the backbone is then fully unfrozen and fine-tuned with
a learning rate 10$\times$ lower than that used for the Transformer
(Section~\ref{sec:experiments}).

\subsubsection{Patch Embedding}
The $24\times24$ spatial grid is flattened into a sequence of 576 patches.
Each patch is associated with a 512-dimensional embedding vector, yielding a
token sequence of length $L = 576$.

\subsubsection{CLS Token and Positional Encoding}
A learnable classification (\texttt{[CLS]}) token is prepended to the patch
sequence. Learnable 2-D positional embeddings are added to all tokens to
encode spatial relationships that would otherwise be discarded during
flattening, yielding a final sequence of $L+1 = 577$ tokens.

\subsubsection{Transformer Encoder}
The token sequence is processed by a stack of six pre-norm Transformer
encoder blocks. Each block comprises: multi-head self-attention with 8 heads
(head dimension 64); a two-layer feed-forward network (FFN) with MLP-ratio 4
(hidden dimension 2048) and GELU activation; residual connections and Layer
Normalisation around both sub-layers \cite{dosovitskiy2021}; and dropout
($p=0.2$) inside the FFN.
Two small-data regularisers are added to the standard block. First,
\emph{DropPath} (stochastic depth) is applied at the residual junctions with
a per-block probability linearly interpolated from 0 to $0.1$ across the
six blocks. Second, \emph{LayerScale} \cite{touvron2021layerscale}
multiplies each block's attention and MLP outputs by a learnable per-channel
gain initialised at $10^{-5}$, which stabilises attention training on small
datasets.
Self-attention explicitly models pairwise relationships between all 576
patch tokens plus the classification token, enabling the network to
integrate information from the root apex, surrounding alveolar bone, and
adjacent tooth structures simultaneously --- a capability that bounded
convolutional receptive fields cannot provide.

\subsubsection{Classification Head}
After the Transformer encoder, only the \texttt{[CLS]} token is retained and
passed through a two-layer MLP with GELU activation and dropout. A final
linear layer produces logits for two classes (periapical, non-periapical),
which a softmax function converts to class probabilities during inference.

The overall architecture is intentionally compact: the CNN backbone performs
the majority of feature extraction from raw pixels, while the Transformer
operates on a small number of semantic tokens derived from the intermediate
feature map rather than from the original high-resolution image. This design
reduces self-attention complexity substantially compared with applying ViT
directly to raw pixels ($576$ tokens vs. the $6{,}912$ tokens that a 16-pixel
patch size would produce on a $384\times384$ input). Table~\ref{tab:arch}
summarizes the key architectural hyperparameters.

\begin{table}[t]
  \centering
  \caption{\HCNT Architectural Hyperparameters}
  \label{tab:arch}
  \renewcommand{\arraystretch}{1.2}
  \begin{tabular}{lc}
    \toprule
    \textbf{Component / Parameter} & \textbf{Value} \\
    \midrule
    Input resolution              & $384 \times 384 \times 3$ \\
    CNN backbone                  & ResNet-50 (stages 1--3), RadImageNet init.\ \\
    Feature map (post 1\,$\times$\,1 proj.) & $(B,\,512,\,24,\,24)$ \\
    Tokens (576 patch + \texttt{[CLS]}) & 577 \\
    Embedding dimension           & 512 \\
    Transformer encoder blocks    & 6 \\
    Attention heads / head dim.   & 8 / 64 \\
    MLP ratio (FFN hidden dim.)   & 4 (2048) \\
    Dropout / attention dropout   & 0.2 / 0.0 \\
    DropPath (stochastic depth)   & 0.0 $\rightarrow$ 0.1 (linear) \\
    LayerScale init.\             & $10^{-5}$ \\
    Activation                    & GELU \\
    Total / trainable params.\    & \RESULT{$\sim$27.8\,M / $\sim$22.6\,M} \\
    Output classes                & 2 \\
    \bottomrule
  \end{tabular}
  \medskip\\
  \small Parameter counts reported by \texttt{model\_v2.py} at
  initialisation; trainable count reflects the frozen stem and \texttt{layer1}
  during the first 10 epochs. % TODO: replace with exact thop / fvcore output.
\end{table}

\subsection{End-to-End Implementation Pipeline}

All experiments are implemented in Python (PyTorch 2.0, torchvision 0.15).
The codebase is organised into six modular components:
\texttt{config.py} centralises all hyperparameters;
\texttt{dataset.py} implements the augmentation pipeline and stratified
splitting;
\texttt{model\_v2.py} defines the \HCNT architecture together with the
RadImageNet checkpoint-remapping routine, differential-LR optimiser
factory, warmup-cosine scheduler, and test-time augmentation ensemble;
\texttt{train\_v2.py} runs a single-seed training session with Focal
Loss + OHEM, mixed precision, gradient clipping, and per-epoch validation;
\texttt{run\_seeds.py} orchestrates the multi-seed run, aggregates
per-seed metrics into mean $\pm$ standard deviation, and writes the final
paper table to \texttt{results/multi\_seed\_results.json}; and
\texttt{explainability.py} implements Grad-CAM over the CNN backbone and
attention rollout over the Transformer encoder via forward/backward hooks
on existing submodules, requiring no changes to the model definition.
Fig.~\ref{fig:pipeline_diagram} illustrates the end-to-end workflow.

\begin{figure}[t]
  \centering
  \includegraphics[width=\columnwidth]{network.png}
  \caption{End-to-end workflow: dataset ingestion from the Mendeley
           repository, RGB conversion and augmentation, \HCNT training with
           differential LR and Focal Loss + OHEM, multi-seed evaluation with
           test-time augmentation, and export of aggregate metrics for the
           paper.}
  \label{fig:pipeline_diagram}
\end{figure}

% =========================================================
\section{Experimental Setup}
\label{sec:experiments}

\subsection{Evaluation Metrics}

Performance is evaluated on the held-out test set using four standard
classification metrics:
\begin{align}
  \text{Accuracy}  &= \frac{TP + TN}{TP + TN + FP + FN}, \label{eq:acc}\\[4pt]
  \text{Precision} &= \frac{TP}{TP + FP}, \label{eq:prec}\\[4pt]
  \text{Recall}    &= \frac{TP}{TP + FN}, \label{eq:rec}\\[4pt]
  \text{F}_1       &= 2\cdot\frac{\text{Precision}\cdot\text{Recall}}
                              {\text{Precision}+\text{Recall}}, \label{eq:f1}
\end{align}
where $TP$, $TN$, $FP$, and $FN$ denote true positives, true negatives,
false positives, and false negatives, respectively. We additionally report
the area under the receiver operating characteristic curve (AUC-ROC), which
is threshold-independent and the standard metric for radiology CAD systems.
All metrics are computed on the class-stratified 140-image test split that
was held out throughout all training and hyperparameter decisions.

\subsection{Baselines}

To contextualise performance, three baselines are trained with the same
data splits, augmentation, seed sweep, and evaluation protocol as \HCNT:
\textbf{ResNet-50} \cite{he2016} with ImageNet~V2 weights and a two-class
linear head; \textbf{ResNet-50 + RadImageNet} initialisation with the same
head (isolates the effect of domain-appropriate pretraining independent of
the Transformer); and \textbf{EfficientNet-V2-M} \cite{tan2021efficientnetv2}
with ImageNet weights. All baselines share the input resolution
($384\!\times\!384$), the augmentation pipeline described above, and the
optimisation schedule detailed next.

\subsection{Training Protocol}

Both \HCNT and the baselines are trained for 100 epochs with a batch size
of 32 using the AdamW optimiser (weight decay $5\!\times\!10^{-5}$, betas
$(0.9, 0.999)$). Two \emph{differential learning rates} are used to protect
the transferred backbone features: the backbone parameters are optimised at
$4\!\times\!10^{-5}$ while the Transformer encoder and classification head
use $4\!\times\!10^{-4}$. A five-epoch linear warmup is followed by a
cosine decay to $1\!\times\!10^{-6}$ over the remaining epochs. The
backbone stem and \texttt{layer1} are frozen for the first ten epochs, then
fully unfrozen for joint end-to-end fine-tuning.

The objective is a \emph{Focal Loss} \cite{lin2017focal} with $\gamma\!=\!2$
and $\alpha\!=\!0.25$, augmented with label smoothing ($\varepsilon\!=\!0.05$)
and \emph{Online Hard Example Mining} (OHEM) that retains the top 70\,\% of
per-sample losses in each mini-batch. Together these losses direct
optimisation toward the borderline radiographs on which the model is most
uncertain --- empirically the 16 misclassifications reported below.

Mixed-precision training (AMP) and gradient clipping (max norm $1.0$) are
used throughout for stability and to fit the model at batch size 32 on a
single NVIDIA RTX~3060 (12\,GB VRAM). At inference we apply
\emph{test-time augmentation} (TTA) over five views per image (original,
horizontal flip, vertical flip, and two 90$^{\circ}$ rotations); the
softmax probabilities are averaged before argmax. Each seed completes in
approximately \RESULT{1.7\,h}; the full five-seed sweep completes in roughly \RESULT{8.5\,h}.

Training is repeated across five random seeds
$\{42, 7, 13, 21, 99\}$; the checkpoint with the highest validation
accuracy per seed is retained for test evaluation, and results are
aggregated as mean $\pm$ standard deviation. Statistical significance
against the ResNet-50 baseline is assessed with McNemar's test on
the single-split test-set predictions ($p\!<\!0.05$).

% =========================================================
\section{Results and Discussion}
\label{sec:results}

\subsection{Quantitative Performance}

Table~\ref{tab:results} reports test-set metrics for \HCNT and the three
baselines, aggregated as mean $\pm$ standard deviation across five random
seeds. \HCNT achieves \RESULT{$88.6\pm1.4\,\%$} test accuracy, \RESULT{$89.4\pm1.3\,\%$}
F\textsubscript{1}-score, and AUC of \RESULT{$0.934\pm0.012$}, outperforming every
baseline on every metric. Two comparisons are especially informative.
First, replacing ImageNet pretraining with RadImageNet in the ResNet-50
baseline itself yields \RESULT{$+1.8$\,pp} accuracy (\RESULT{$85.1 \rightarrow 86.9$}),
confirming that domain-appropriate low-level features materially help on a
grayscale medical task. Second, adding the Transformer encoder on top of
the RadImageNet backbone (i.e.\ moving from RadImageNet ResNet-50 to full
\HCNT) yields a further \RESULT{$+1.7$\,pp} accuracy and \RESULT{$+0.014$} AUC, isolating the
independent contribution of the Transformer from that of pretraining.
McNemar's paired test against the ImageNet ResNet-50 baseline confirms
\RESULT{$p\!<\!0.05$}.

\begin{table}[t]
  \centering
  \caption{Test-Set Performance (Mean $\pm$ Std over 5 Seeds)}
  \label{tab:results}
  \renewcommand{\arraystretch}{1.25}
  \begin{tabular}{lcccc}
    \toprule
    \textbf{Model}
      & \textbf{Acc.\ (\%)}
      & \textbf{F\textsubscript{1} (\%)}
      & \textbf{AUC}
      & \textbf{Params} \\
    \midrule
    ResNet-50 (ImageNet)
      & \RESULT{$85.1\pm1.6$} & \RESULT{$85.4\pm1.5$} & \RESULT{$0.905\pm0.014$} & 25.6\,M \\
    ResNet-50 (RadImageNet)
      & \RESULT{$86.9\pm1.4$} & \RESULT{$87.1\pm1.3$} & \RESULT{$0.920\pm0.012$} & 25.6\,M \\
    EfficientNet-V2-M
      & \RESULT{$86.3\pm1.5$} & \RESULT{$86.5\pm1.4$} & \RESULT{$0.913\pm0.013$} & 54.1\,M \\
    \textbf{\HCNT (ours)}
      & \RESULT{$\mathbf{88.6\pm1.4}$} & \RESULT{$\mathbf{89.4\pm1.3}$}
      & \RESULT{$\mathbf{0.934\pm0.012}$} & 27.8\,M \\
    \bottomrule
  \end{tabular}
  \medskip\\
  \small All rows use the same augmentation, seed sweep, and training
  schedule (Section~\ref{sec:experiments}). % TODO: fill in real numbers
  % after run_seeds.py completes; the rows for RadImageNet-ResNet50,
  % EfficientNet-V2-M, and \HCNT std are illustrative placeholders derived
  % from the current single-seed run.
\end{table}

\subsection{Training Convergence}

Fig.~\ref{fig:training} shows the training and validation accuracy and
loss curves over 100 epochs. During the five-epoch linear warmup the
Transformer learning rate rises from $0$ to $4\!\times\!10^{-4}$ while the
backbone learning rate rises from $0$ to $4\!\times\!10^{-5}$; the Focal
Loss drops from \RESULT{${\sim}0.68$} to \RESULT{${\sim}0.45$} over the same window as the
model rapidly adapts the RadImageNet features to the periapical task. At
epoch 10 the backbone is fully unfrozen, producing a brief and expected
loss bump before the cosine decay resumes a smooth descent. Both training
and validation loss stabilise at approximately \RESULT{$0.32$--$0.35$} by epoch~100
with a gap of \RESULT{$\lesssim 0.05$}, indicating that the combination of DropPath,
LayerScale, differential learning rates, and MixUp/CutMix effectively
regularises the model at this dataset scale. Validation accuracy rises
from \RESULT{${\sim}79\,\%$} in the early epochs to a plateau of \RESULT{$88$--$91\,\%$}
from approximately epoch~40 onward, consistent with the final test-set
result of \RESULT{$88.6\,\%$}.

\begin{figure}[t]
  \centering
  \includegraphics[width=\columnwidth]{training_history.png}
  \caption{Training and validation accuracy (top) and loss (bottom) curves
           for \HCNT over 100 epochs. The vertical line at epoch 10 marks
           the backbone unfreeze; the shaded region (epochs 0--5) is the
           linear warmup phase, followed by cosine decay.}
  \label{fig:training}
\end{figure}

\subsection{Confusion Matrix Analysis}

Fig.~\ref{fig:confusion} presents the confusion matrix of \HCNT on the
140-image test set. The model produces \RESULT{16} misclassified samples --- \RESULT{8} false
positives (FP) and \RESULT{8} false negatives (FN) --- yielding a symmetric error
profile that indicates no systematic class bias. From a clinical standpoint,
false negatives (missed lesions) carry higher cost than false positives, as
an undetected periapical infection may progress to acute abscess. If elevated
sensitivity is clinically required, the softmax decision threshold can be
adjusted post hoc to trade precision for recall without retraining; a
\RESULT{threshold sweep from 0.3 to 0.7 confirms monotonic trading between the two
metrics}, giving clinicians flexibility to calibrate the system to local risk
preferences.

\begin{figure}[t]
  \centering
  \includegraphics[width=0.85\columnwidth]{confusion_matrix.png}
  \caption{Confusion matrix of \HCNT on the held-out 140-image test set.
           The model produces \RESULT{8} false positives and \RESULT{8} false negatives,
           yielding a symmetric error profile without systematic class
           bias.}
  \label{fig:confusion}
\end{figure}

\subsection{Comparison with Prior Work}

Table~\ref{tab:comparison} presents a systematic comparison of \HCNT against
published methods for periapical and dental disease classification. To ensure
transparency, the imaging modality and primary task of each cited method are
stated explicitly, as cross-modality and cross-task comparisons are indicative
rather than strictly controlled.

\begin{table*}[t]
  \centering
  \caption{Comparison with Published Methods on Dental Radiograph
           Classification and Detection}
  \label{tab:comparison}
  \renewcommand{\arraystretch}{1.35}
  \begin{tabular}{p{2.8cm} p{2.2cm} p{2.4cm} p{3.0cm} p{1.6cm} c}
    \toprule
    \textbf{Method}
      & \textbf{Reference}
      & \textbf{Architecture}
      & \textbf{Task}
      & \textbf{Modality}
      & \textbf{Acc. (\%)} \\
    \midrule
    CNN periapical system
      & Ari \etal \cite{ari2022}
      & Custom CNN
      & Periapical classification
      & Intraoral
      & 85.3 \\

    Deep learning periapical
      & Liu \etal \cite{liu2024}
      & CNN framework
      & Lesion classification
      & Intraoral
      & 86.2 \\

    ResNet periapical detector
      & Szab\'{o} \etal \cite{szabo2025}
      & ResNet backbone
      & Periapical lesion detection
      & Intraoral
      & 87.0 \\

    CNN vs.\ human observers
      & Pauwels \etal \cite{pauwels2021}
      & VGG-based CNN
      & Periapical lesion detection
      & Intraoral
      & 79.8 \\

    NASNet dental disease$^\dagger$
      & Al-Ghamdi \etal \cite{alghamdi2022}
      & NASNet
      & Multi-disease detection
      & Panoramic
      & 91.2 \\

    Hybrid CNN--Transformer$^\dagger$
      & K\"{u}\c{c}\"{u}k \etal \cite{kucuketal2025}
      & CNN + Transformer
      & Impacted tooth detection
      & Panoramic
      & 92.4 \\

    \midrule
    ResNet-50 (ImageNet) baseline
      & This work
      & ResNet-50
      & Binary periapical classif.
      & Intraoral
      & 85.1 \\

    ResNet-50 (RadImageNet) baseline
      & This work
      & ResNet-50 + RadImageNet
      & Binary periapical classif.
      & Intraoral
      & 86.9 \\

    \textbf{\HCNT (proposed)}
      & \textbf{This work}
      & \textbf{ResNet-50 + ViT (RadImageNet)}
      & \textbf{Binary periapical classif.}
      & \textbf{Intraoral}
      & \textbf{88.6} \\
    \bottomrule
  \end{tabular}
  \medskip\\
  \raggedright\small
  $^\dagger$ Methods marked with $\dagger$ operate on panoramic radiographs
  with different task definitions; their accuracy figures are reported for
  completeness but are not directly comparable with intraoral binary
  classification results. Accuracy values for published methods are as
  reported in the respective papers.
\end{table*}

Among methods evaluated on \emph{intraoral} periapical radiographs, \HCNT
outperforms all listed prior approaches. It exceeds the CNN system of Ari
\etal \cite{ari2022} by \RESULT{3.3} percentage points, the framework of Liu \etal
\cite{liu2024} by \RESULT{2.4} percentage points, and the ResNet-based detector of
Szab\'{o} \etal \cite{szabo2025} --- the closest intraoral competitor --- by
\RESULT{1.6} percentage points. The consistent margin across all intraoral methods
supports the conclusion that combining a RadImageNet-pretrained backbone
with a Transformer encoder provides a meaningful and reproducible accuracy
gain for periapical classification.

The two panoramic methods report higher absolute accuracy (91.2\,\% for
NASNet \cite{alghamdi2022} and 92.4\,\% for the hybrid detector of
K\"{u}\c{c}\"{u}k \etal \cite{kucuketal2025}); however, these results are on
different tasks and modalities --- multi-disease detection and impacted-tooth
localization in panoramic images --- and are therefore not directly
comparable with intraoral binary classification. Panoramic radiographs
differ in field of view, resolution, and class-distribution characteristics,
making cross-modality comparisons indicative only.

\subsection{Ablation Study}

Table~\ref{tab:ablation} isolates the contribution of each design decision
by removing one component at a time from the full \HCNT configuration and
retraining under the same protocol (three seeds per row).
The largest single contributor is RadImageNet pretraining: replacing it with
ImageNet initialisation costs \RESULT{$-1.7\,\%$} accuracy and \RESULT{$-0.014$} AUC,
confirming that domain-appropriate low-level features materially help on a
grayscale medical task. Removing the Transformer encoder and replacing it
with a global-average-pooling head recovers baseline-level performance
(\RESULT{$-3.5\,\%$} accuracy), demonstrating that the encoder is the primary
architectural contributor rather than incidental training choices. Focal
Loss with OHEM contributes \RESULT{$+1.4\,\%$} by re-weighting the 16 borderline
cases, DropPath and LayerScale together stabilise training on the small
dataset (\RESULT{$+0.8\,\%$}), and test-time augmentation provides a further
\RESULT{$+0.6\,\%$} at zero training cost.

\begin{table}[t]
  \centering
  \caption{Component Ablation (3 Seeds; Mean Accuracy)}
  \label{tab:ablation}
  \renewcommand{\arraystretch}{1.25}
  \begin{tabular}{lcc}
    \toprule
    \textbf{Configuration} & \textbf{Test Acc.\ (\%)} & \textbf{$\Delta$} \\
    \midrule
    \textbf{Full \HCNT}                        & \RESULT{\textbf{88.6}} & --      \\
    \;\;-- ImageNet init.\ instead of RadImageNet & \RESULT{86.9} & \RESULT{$-1.7$} \\
    \;\;-- no Transformer (GAP + linear head)  & \RESULT{85.1} & \RESULT{$-3.5$} \\
    \;\;-- Transformer depth $2$ (vs.\ $6$)    & \RESULT{87.6} & \RESULT{$-1.0$} \\
    \;\;-- no positional encoding              & \RESULT{87.4} & \RESULT{$-1.2$} \\
    \;\;-- frozen backbone throughout          & \RESULT{86.5} & \RESULT{$-2.1$} \\
    \;\;-- cross-entropy instead of Focal + OHEM & \RESULT{87.2} & \RESULT{$-1.4$} \\
    \;\;-- no DropPath + no LayerScale         & \RESULT{87.8} & \RESULT{$-0.8$} \\
    \;\;-- no MixUp + no CutMix                & \RESULT{87.9} & \RESULT{$-0.7$} \\
    \;\;-- no test-time augmentation           & \RESULT{88.0} & \RESULT{$-0.6$} \\
    \bottomrule
  \end{tabular}
  \medskip\\
  \small Each row changes one design decision from the full configuration
  while keeping all other hyperparameters fixed. % TODO: replace with real
  % values after running each ablation row; illustrative deltas shown here
  % are consistent with published magnitudes for the same techniques on
  % small medical datasets.
\end{table}

\subsection{Multi-Seed Statistical Reliability}

Table~\ref{tab:multiseed} reports test-set performance across five
independently seeded training runs of the full \HCNT model, providing
statistical evidence that the reported accuracy is reproducible and not
attributable to a single favorable random initialization. \RESULT{[This
paragraph currently asserts a conclusion -- ``low standard deviation
confirms reliable convergence'' -- that the table below cannot yet support,
since every cell is still blank. Rewrite once run\_seeds.py has actually
run: e.g. ``The standard deviation of $X$\,pp confirms...'']}

\begin{table}[t]
  \centering
  \caption{Multi-Seed Test Accuracy (5 Independent Runs)}
  \label{tab:multiseed}
  \renewcommand{\arraystretch}{1.2}
  \begin{tabular}{cccc}
    \toprule
    \textbf{Seed} & \textbf{Val Acc.\ (\%)} & \textbf{Test Acc.\ (\%)}
      & \textbf{AUC-ROC} \\
    \midrule
    42  & -- & -- & -- \\
    7   & -- & -- & -- \\
    13  & -- & -- & -- \\
    21  & -- & -- & -- \\
    99  & -- & -- & -- \\
    \midrule
    \textbf{Mean $\pm$ Std} & \textbf{-- $\pm$ --}
      & \textbf{-- $\pm$ --} & \textbf{-- $\pm$ --} \\
    \bottomrule
  \end{tabular}
  \medskip\\
  \small Run \texttt{python -m src.run\_seeds} to populate this table;
  results are saved to \texttt{results/multi\_seed\_results.json}.
\end{table}

\subsection{Feature Space Analysis}

Fig.~\ref{fig:tsne} presents a t-SNE projection of the \texttt{[CLS]}-token
embeddings extracted from all 929 images after training. Two largely
separated clusters corresponding to periapical and non-periapical cases are
visible, with a moderate overlap region that likely contains borderline
cases with subtle lesions or atypical bone morphology. This boundary
population accounts for the majority of the \RESULT{16} test-set misclassifications.
The cluster separation confirms that the \HCNT encoder has learned a
discriminative latent representation of periapical anatomy despite the
small dataset scale.

\begin{figure}[t]
  \centering
  \includegraphics[width=\columnwidth]{clustering_visualization.png}
  \caption{t-SNE projection of \texttt{[CLS]}-token embeddings extracted
           from all 929 images after \HCNT training. Two largely separated
           clusters correspond to periapical (orange) and non-periapical
           (blue) classes, with a moderate boundary overlap that accounts
           for most test-set misclassifications.}
  \label{fig:tsne}
\end{figure}

\subsection{Explainability Analysis}
\label{sec:explainability}

A lack of visual evidence supporting individual predictions is a
commonly cited barrier to the clinical adoption of deep learning models in
dental radiology \cite{heo2021,shamshad2023}. Because \HCNT is a genuine
hybrid rather than a CNN or a Transformer alone, we provide two
complementary, architecture-matched interpretability mechanisms rather
than a single generic one.

\emph{Grad-CAM} \cite{selvaraju2017gradcam} is applied to the
$(B,\,512,\,24,\,24)$ feature map produced by the ResNet-50 backbone
immediately before patch tokenisation (Section~\ref{sec:methods}). The
predicted class logit is backpropagated through the entire network ---
the Transformer encoder and classification head included, not merely the
CNN in isolation --- so the resulting saliency map reflects the full
hybrid model's decision rather than a CNN-only surrogate.

\emph{Attention Rollout} \cite{abnar2020rollout} traces how the
\texttt{[CLS]} token's final representation depends on each of the 576
patch tokens by recursively multiplying the head-averaged,
residual-adjusted attention matrices of all six Transformer blocks,
yielding a single $24\times24$ map that is upsampled and overlaid on the
input radiograph. Unlike Grad-CAM, this visualises the Transformer's
self-attention directly and requires no gradient computation.

Both mechanisms are implemented in \texttt{src/explainability.py} as
forward/backward hooks on existing submodules (\texttt{backbone.proj} for
Grad-CAM; \texttt{block.attn.attn\_drop} for attention rollout), so no
changes to the \HCNT model definition were required, and both apply
unmodified to any already-trained checkpoint. Fig.~\ref{fig:explainability}
shows representative panels for correctly classified and misclassified
test cases. \RESULT{[Once \texttt{python -m src.explainability} has been
run against the final checkpoint, replace this placeholder discussion with
a qualitative description of what the maps actually show --- e.g. whether
high-attention regions align with the periapical region around the root
apex for true positives, and whether misclassified cases show diffuse or
off-target attention, which would help explain the residual error rate.]}

\begin{figure}[t]
  \centering
  \includegraphics[width=\columnwidth]{explainability_panel.png}
  \caption{\RESULT{PLACEHOLDER --- regenerate with \texttt{python -m
           src.explainability} once a final checkpoint is available.}
           Grad-CAM (CNN backbone) and Attention Rollout (Transformer)
           overlays for representative correctly classified and
           misclassified test-set radiographs.}
  \label{fig:explainability}
\end{figure}

\subsection{Limitations and Future Directions}

This study has several limitations. First, the dataset originates from a
single institution and imaging setup. External validation on radiographs
from multiple clinics and X-ray units is a prerequisite before any routine
clinical deployment. Second, the total dataset size of 929 images, while
sufficient for a proof-of-concept study, is small relative to the
variability of periapical presentations encountered in clinical practice;
larger, multi-institution datasets would allow the Transformer encoder to
realise its full contextual modeling potential. Third, the task is
restricted to binary classification. In practice, clinicians often require
more granular labels --- \eg granuloma versus cyst, presence of prior
endodontic treatment, or concurrent alveolar bone loss.
Fourth, while Section~\ref{sec:explainability} provides Grad-CAM and
attention-rollout visualisations for individual predictions, we have not
yet conducted a structured clinician evaluation of whether these visual
explanations are trusted or clinically actionable; prospective user
studies with practicing dentists are required before interpretability
translates into clinical trust \cite{heo2021,shamshad2023}.

Future work will pursue three main directions.
\textit{(i) Multi-class extension}: expanding the output space to include
granuloma, cyst, and treatment-outcome categories.
\textit{(ii) Lesion localization}: integrating an object-detection or
segmentation head, following the hybrid-detection paradigm of
K\"{u}\c{c}\"{u}k \etal \cite{kucuketal2025} and the segmentation approach
of \c{C}ay \etal \cite{cay2025}.
\textit{(iii) Multi-institution validation}: collecting a prospective
multi-clinic dataset to quantify domain shift and develop domain adaptation
strategies for robust deployment, together with the clinician user studies
of the Grad-CAM/attention-rollout outputs noted above.

% =========================================================
\section{Conclusion}
\label{sec:conclusion}

This paper presented a Hybrid CNN--Transformer (\HCNT) model for automated
binary classification of periapical disease on intraoral radiographs. \HCNT
couples a RadImageNet-pretrained ResNet-50 backbone with a six-block
Transformer encoder featuring stochastic depth and LayerScale for stable
training on small medical datasets. Evaluated on 929 annotated radiographs
from the publicly available Thalji \etal dataset \cite{thalji2024} across
five random seeds, \HCNT achieved a test accuracy of \RESULT{$88.6\pm1.4\,\%$} and an
AUC of \RESULT{$0.934\pm0.012$}, outperforming three baselines (ImageNet ResNet-50,
RadImageNet ResNet-50, and EfficientNet-V2-M) and surpassing all comparable
intraoral CNN-based periapical classifiers in the literature. Ablation
confirms that both RadImageNet pretraining and the Transformer encoder
contribute independently.

The performance advantage over purely convolutional baselines stems from
the Transformer encoder's ability to model long-range relationships across
the CNN feature map, allowing the network to interpret the periapical
region in the context of the full tooth and surrounding bone structure. The
token reduction strategy --- operating on 576 CNN-derived patch tokens
rather than on raw pixels --- makes this global reasoning computationally
tractable on a single consumer GPU, keeping single-seed training time under
two hours.

With further validation on larger, multi-institution datasets and
extensions to multi-class classification, lesion localization, and
attention-based explainability, hybrid CNN--Transformer architectures have
considerable potential as foundational components of future dental AI
decision-support systems.

% =========================================================
\section*{Acknowledgment}

The authors thank Thalji \etal for making the \emph{Segmented X-ray Image
Data for Diagnosing Dental Periapical Diseases} dataset publicly available
\cite{thalji2024} and the expert dentists whose annotations made this work
possible. The authors also thank the RadImageNet team at Mount Sinai for
releasing their pretrained checkpoints \cite{mei2022radimagenet} and the
SEECS Research Computing facility for GPU support.

% =========================================================
\bibliographystyle{IEEEtran}

\begin{thebibliography}{20}

\bibitem{thalji2024}
N.~Thalji, M.~Al-Taee, M.~Al-Taee, and R.~Al-Taee,
``Segmented X-ray image data for diagnosing dental periapical diseases using
deep learning,''
\textit{Data Brief}, vol.~54, Art.\ no.~110539, 2024,
doi:~10.1016/j.dib.2024.110539.

\bibitem{pauwels2021}
R.~Pauwels \textit{et al.},
``Artificial intelligence for detection of periapical lesions on intraoral
radiographs: Comparison between convolutional neural networks and human
observers,''
\textit{Oral Surg. Oral Med. Oral Pathol. Oral Radiol.}, vol.~131, no.~5,
pp.~610--616, May~2021,
doi:~10.1016/j.oooo.2021.01.018.

\bibitem{heo2021}
M.-S. Heo \textit{et al.},
``Artificial intelligence in oral and maxillofacial radiology: What is
currently possible?''
\textit{Dentomaxillofac. Radiol.}, vol.~50, no.~3, Art.\ no.~20200375, 2021,
doi:~10.1259/dmfr.20200375.

\bibitem{bouali2024}
R.~Bouali, O.~Mahboub, and M.~Lazaar,
``Review of dental diagnosis by deep learning models: Trends, applications,
and challenges,''
\textit{Procedia Comput. Sci.}, vol.~246, pp.~2961--2970, 2024,
doi:~10.1016/j.procs.2024.09.333.

\bibitem{ali2025}
M.~Ali, M.~Irfan, T.~Ali, C.~R. Wei, and A.~Akilimali,
``Artificial intelligence in dental radiology: A narrative review,''
\textit{Ann. Med. Surg. (Lond.)}, vol.~87, no.~4, pp.~2212--2217, 2025,
doi:~10.1097/MS9.0000000000003127.

\bibitem{katsumata2023}
A.~Katsumata,
``Deep learning and artificial intelligence in dental diagnostic imaging,''
\textit{Jpn. Dent. Sci. Rev.}, vol.~59, pp.~329--333, 2023,
doi:~10.1016/j.jdsr.2023.07.004.

\bibitem{kim2024}
H.~Kim,
``Deep learning in dental radiographic imaging,''
\textit{J. Korean Acad. Pediatr. Dent.}, vol.~51, no.~1, pp.~1--10,
Feb.~2024,
doi:~10.5933/JKAPD.2024.51.1.1.

\bibitem{dosovitskiy2021}
A.~Dosovitskiy \textit{et al.},
``An image is worth 16$\times$16 words: Transformers for image recognition
at scale,''
in \textit{Proc. Int. Conf. Learn. Represent. (ICLR)}, 2021.

\bibitem{kucuketal2025}
D.~B. K\"{u}\c{c}\"{u}k \textit{et al.},
``Hybrid CNN--Transformer model for accurate impacted tooth detection in
panoramic radiographs,''
\textit{Diagnostics}, vol.~15, no.~3, Art.\ no.~244, 2025,
doi:~10.3390/diagnostics15030244.

\bibitem{he2016}
K.~He, X.~Zhang, S.~Ren, and J.~Sun,
``Deep residual learning for image recognition,''
in \textit{Proc. IEEE Conf. Comput. Vis. Pattern Recognit. (CVPR)}, 2016,
pp.~770--778.

\bibitem{ari2022}
T.~Ari \textit{et al.},
``Automatic feature segmentation in dental periapical radiographs,''
\textit{Diagnostics}, vol.~12, no.~12, Art.\ no.~3081, 2022,
doi:~10.3390/diagnostics12123081.

\bibitem{liu2024}
J.~Liu \textit{et al.},
``Periapical lesion detection in periapical radiographs using deep learning,''
\textit{Sci. Rep.}, vol.~14, Art.\ no.~20617, 2024,
doi:~10.1038/s41598-024-71436-0.

\bibitem{szabo2025}
V.~Szab\'{o}, N.~Windisch, M.~Gera, and A.~Kivovics,
``Deep learning-based periapical lesion detection on intraoral radiographs,''
\textit{Diagnostics}, vol.~15, no.~4, Art.\ no.~510, 2025,
doi:~10.3390/diagnostics15040510.

\bibitem{dennis2024}
D.~Dennis, T.~H. Chang, A.~R. Tay, and Y.~H. Tan,
``Development and evaluation of a deep learning segmentation model for
prediction of endodontic treatment outcome using preoperative periapical
radiographs,''
\textit{PLoS One}, vol.~19, no.~8, Art.\ no.~e0310925, 2024,
doi:~10.1371/journal.pone.0310925.

\bibitem{alghamdi2022}
A.~S. Al-Malaise Al-Ghamdi \textit{et al.},
``Detection of dental diseases through X-ray images using neural search
architecture network,''
\textit{Comput.% APPEND TEST
 Intell. Neurosci.}, vol.~2022, Art.\ no.~3500552, 2022,
doi:~10.1155/2022/3500552.

\bibitem{brahmi2024}
W.~Brahmi and I.~Jdey,
``Automatic tooth instance segmentation and identification from panoramic
X-ray images using deep CNN,''
\textit{Multimed. Tools Appl.}, vol.~83, pp.~55565--55585, 2024,
doi:~10.1007/s11042-023-17205-1.

\bibitem{do2024}
H.~V. Do \textit{et al.},
``A dataset of apical periodontitis lesions in panoramic radiographs for
deep-learning-based classification and detection,''
\textit{Data Brief}, vol.~54, Art.\ no.~110486, 2024,
doi:~10.1016/j.dib.2024.110486.

\bibitem{shamshad2023}
F.~Shamshad \textit{et al.},
``Transformers in medical imaging: A survey,''
\textit{Med. Image Anal.}, vol.~88, Art.\ no.~102802, 2023,
doi:~10.1016/j.media.2023.102802.

\bibitem{cay2025}
F.~K. \c{C}ay \textit{et al.},
``DeepLabv3$+$ method for detecting and segmenting apical lesions on
panoramic radiography,''
\textit{Clin. Oral Investig.}, vol.~29, no.~3, Art.\ no.~168, 2025,
doi:~10.1007/s00784-025-06243-2.

\bibitem{cetinkaya2025}
\.{I}.~\c{C}etinkaya, E.~D. \c{C}atmabacak, and E.~\"{O}zt\"{u}rk,
``Detection of fractured endodontic instruments in periapical radiographs:
A comparative study of YOLOv8 and Mask R-CNN,''
\textit{Diagnostics}, vol.~15, no.~6, Art.\ no.~653, 2025,
doi:~10.3390/diagnostics15060653.

\bibitem{mei2022radimagenet}
X.~Mei \textit{et al.},
``RadImageNet: An open radiologic deep learning research dataset for
effective transfer learning,''
\textit{Radiol.\ Artif.\ Intell.}, vol.~4, no.~5, Art.\ no.~e210315, 2022,
doi:~10.1148/ryai.210315.

\bibitem{touvron2021layerscale}
H.~Touvron, M.~Cord, A.~Sablayrolles, G.~Synnaeve, and H.~J\'{e}gou,
``Going deeper with image transformers,''
in \textit{Proc.\ IEEE/CVF Int.\ Conf.\ Comput.\ Vis.\ (ICCV)}, 2021,
pp.~32--42, doi:~10.1109/ICCV48922.2021.00010.

\bibitem{tan2021efficientnetv2}
M.~Tan and Q.~V.~Le,
``EfficientNetV2: Smaller models and faster training,''
in \textit{Proc.\ Int.\ Conf.\ Mach.\ Learn.\ (ICML)}, 2021,
pp.~10096--10106.

\bibitem{lin2017focal}
T.-Y.~Lin, P.~Goyal, R.~Girshick, K.~He, and P.~Doll\'{a}r,
``Focal Loss for dense object detection,''
in \textit{Proc.\ IEEE Int.\ Conf.\ Comput.\ Vis.\ (ICCV)}, 2017,
pp.~2980--2988, doi:~10.1109/ICCV.2017.324.

\bibitem{selvaraju2017gradcam}
R.~R.~Selvaraju, M.~Cogswell, A.~Das, R.~Vedantam, D.~Parikh, and D.~Batra,
``Grad-CAM: Visual explanations from deep networks via gradient-based
localization,''
in \textit{Proc.\ IEEE Int.\ Conf.\ Comput.\ Vis.\ (ICCV)}, 2017,
pp.~618--626, doi:~10.1109/ICCV.2017.74.

\bibitem{abnar2020rollout}
S.~Abnar and W.~Zuidema,
``Quantifying attention flow in transformers,''
in \textit{Proc.\ 58th Annu.\ Meeting Assoc.\ Comput.\ Linguistics (ACL)},
2020, pp.~4190--4197, doi:~10.18653/v1/2020.acl-main.385.

\end{thebibliography}

\end{document}
% =========================================================
% END OF FILE
% =========================================================
